% Generated draft: requires mathematical and source review.
\documentclass[11pt]{article}
\usepackage[margin=25mm]{geometry}
\usepackage{fontspec}
\setmainfont{Latin Modern Roman}
\usepackage{amsmath,amssymb,mathtools,unicode-math}
\setmathfont{Latin Modern Math}
\usepackage{enumitem,needspace,xurl,hyperref}
\setlength{\parindent}{0pt}
\setlength{\parskip}{5pt}
\setlength{\emergencystretch}{3em}
\newcommand{\N}{\mathbb{N}}
\newcommand{\Z}{\mathbb{Z}}
\newcommand{\Q}{\mathbb{Q}}
\newcommand{\R}{\mathbb{R}}
\newcommand{\C}{\mathbb{C}}
\newcommand{\F}{\mathbb{F}}
\newcommand{\A}{\mathbb{A}}
\newcommand{\PP}{\mathbb{P}}
\newcommand{\E}{\mathbb{E}}
\newcommand{\eps}{\varepsilon}
\newcommand{\dd}{\,\mathrm{d}}
\begin{document}
\section*{UnsolvedMath ID: 10400144}
\textbf{Source problem code:} catalogue entry
\subsection{Definitions and mathematical statement}
An integral homology three-sphere is a closed connected three-manifold \(M\) with \(H_*(M;\Z)\cong H_*(S^3;\Z)\); the orientation convention needed for the quantum invariant is to be retained. For each parameter \(\mathfrak g\) ``as above'' in the source, let \(R_1\), \(d\), and \(\tau_\zeta^{\mathfrak g}(M)\) be the ring, positive integer and quantum invariant specified by the missing preceding text. Neither the admissible class of \(\mathfrak g\) nor these definitions is reproduced in the excerpt.

The conjecture attributed to K. Habiro and T. Le asserts that there exists a unique invariant \(I^{\mathfrak g}\) assigning to every integral homology three-sphere \(M\) an element
\[I^{\mathfrak g}(M)\in R_1\]
such that for every root of unity \(\zeta\) of exact order \(r\) divisible by \(d\),
\[\operatorname{ev}_\zeta\bigl(I^{\mathfrak g}(M)\bigr)
=\left.I^{\mathfrak g}(M)\right|_{q=\zeta}
=\tau_\zeta^{\mathfrak g}(M).\]
Here \(\operatorname{ev}_\zeta\) denotes the source's specialization at \(q=\zeta\), whose domain and target must be defined along with \(R_1\). Pointwise uniqueness means that for each fixed \(\mathfrak g,M\) there is exactly one element of \(R_1\) satisfying all these equalities simultaneously. Equivalently, the requested invariant is a lift of the compatible family of root-of-unity values to the specified ring, unique there. Uniqueness depends on the evaluation family, not merely on existence of the individual quantum values. The statement does not restrict to prime-order roots or permit a separate lift for each root. The summary relates the conjecture to the Habiro ring, integral quantum invariants and cyclotomic expansions, but supplies no identification of \(R_1\) with a particular completion; none is imposed here.
\subsection{Short English statement}
Find a unique ring-valued invariant of each integral homology three-sphere whose evaluations at every allowed root of unity recover the corresponding quantum invariants. The ring and allowed Lie-theoretic data are defined only in missing preceding text.
\subsection{Sources}
\begin{itemize}
\item Local UnsolvedMath registry, permanent ID \(10400144\), exact source code \texttt{catalogue entry}.
\item K. Habiro and T. Le, attribution in Ohtsuki, \emph{Problems on invariants of knots and 3-manifolds} (2002), Conjecture 7.29, PDF page 119: \url{https://msp.org/gtm/2002/04/gtm-2002-04-024s.pdf}. Registry access 2026-07-29, pdf-font-structure extraction.
\item The research summary mentions Habiro's cyclotomic-expansion work without a specific bibliographic citation.
\end{itemize}
\end{document}
